\section*{Chapter \ref{ch:b3:computational-chemistry} --- Computational Chemistry: Hartree--Fock and DFT}

\begin{solution}{exo:b3:computational-chemistry:1}
$\dd E/\dd\alpha = \frac32 - \sqrt{2/\pi\alpha} = 0$ gives $\alpha = 8/9\pi = 0.283\,
a_0^{-2}$ and $E = \frac32\alpha - 2\sqrt{2\alpha/\pi} = -4/3\pi = \qty{-0.4244}{\hartree}$,
15\,\% above the exact $-0.5$: a single Gaussian has neither the cusp nor the
tail of the $1s$ function.
\end{solution}

\begin{solution}{exo:b3:computational-chemistry:2}
STO-3G: 5 functions per C ($1s$, $2s$, three $2p$) and 1 per H: $6 \times 5 + 6 =
36$. 6-31G(d): per C, 1 (core) + $2 \times 4$ (split valence) + 6 ($d$) $= 15$;
per H, 2: $6 \times 15 + 6 \times 2 = 102$.
\end{solution}

\begin{solution}{exo:b3:computational-chemistry:3}
The electronic Hamiltonian does not contain the nuclear masses: within the
\omterm{def:b3:computational-chemistry:born-oppenheimer}{Born--Oppenheimer approximation} the PES, hence $r_e$ and $k = U''(r_e)$, is the
same. $\tilde\omega = \sqrt{k/\mu}/2\pi c$ depends on the \omterm{def:b3:quantum-model-systems:oscillator}{reduced mass}, which
doubles: ratio $\sqrt2 = 1.414$ (measured 1.413).
\end{solution}

\begin{solution}{exo:b3:computational-chemistry:4}
Eight atoms give $3N - 6 = 18$ vibrations: 17 real and one imaginary. One
negative Hessian \omterm{def:b3:quantum-model-systems:operator}{eigenvalue}: a first-order saddle point, a transition
structure (here, for example, of a rotation or an elimination).
\end{solution}

\begin{solution}{exo:b3:computational-chemistry:5}
$\zeta = 27/16$, $E = -(27/16)^2 = \qty{-2.8477}{\hartree} = \qty{-77.49}{eV}$.
Error $-2.8477 + 2.9034 = \qty{0.0557}{\hartree} = \qty{1.52}{eV}$, that is
1.9\,\% of the total energy --- but larger than many reaction energies.
\end{solution}

\begin{solution}{exo:b3:computational-chemistry:6}
$(-1 - E)(-0.5 - E) - (-0.2 - 0.3E)^2 = 0$, i.e. $0.91E^2 + 1.38E + 0.46 = 0$:
$E = \qty{-1.0217}{\hartree}$ and \qty{-0.4947}{\hartree}. Yes: mixing in the
second function lowers the energy below $H_{11}$, as the \omterm{thm:b3:computational-chemistry:variational}{variational principle}
allows.
\end{solution}

\begin{solution}{exo:b3:computational-chemistry:7}
$(102/36)^4 = 64$: about 640 minutes, some 11 hours.
\end{solution}

\begin{solution}{exo:b3:computational-chemistry:8}
$0.5782 \times 27.211 = \qty{15.73}{eV}$, 0.30 eV above the measured value. Frozen
orbitals: the real ion relaxes and is lower, so Koopmans' value is too high.
Missing correlation: the neutral molecule (two electrons) has more \omterm{def:b3:computational-chemistry:correlation}{correlation
energy} than the ion (one electron), which makes the true value larger.
Here the first error wins.
\end{solution}

\begin{solution}{exo:b3:computational-chemistry:9}
$-0.5960 - 2(-0.4666) = \qty{0.3372}{\hartree} = \qty{9.18}{eV}$ above two atoms.
The RHF function keeps 50\,\% of \ce{H+ H-}, and separating a proton and a
hydride costs the difference between the ionisation energy of H and the
electron affinity of H, which the energy averages in.
\end{solution}

\begin{solution}{exo:b3:computational-chemistry:10}
$E(c_1^2 + 2c_1c_2S + c_2^2) = c_1^2H_{11} + 2c_1c_2H_{12} + c_2^2H_{22}$.
Differentiating with respect to $c_1$ with $\partial E/\partial c_1 = 0$:
$E(2c_1 + 2c_2S) = 2c_1H_{11} + 2c_2H_{12}$, i.e. $(H_{11} - E)c_1 + (H_{12} -
ES)c_2 = 0$; similarly $(H_{12} - ES)c_1 + (H_{22} - E)c_2 = 0$.
\end{solution}

\begin{solution}{exo:b3:computational-chemistry:11}
$k \approx [E(1.30) - 2E(1.35) + E(1.40)]/(0.05)^2 = 0.001417/0.0025 =
\qty{0.567}{\hartree\per\bohr\squared} = \qty{882}{N/m}$. With $\mu = m_H/2 =
\qty{8.37e-28}{kg}$, $\omega = \sqrt{k/\mu} = \qty{1.03e15}{s^{-1}}$ and
$\tilde\omega = \omega/2\pi c = \qty{5452}{cm^{-1}}$, 24\,\% above the measured
$\tilde\omega_e$: a minimal-basis RHF curve is too steep (no correlation, too
small a basis), which is why computed frequencies are scaled down.
\end{solution}

\begin{solution}{exo:b3:computational-chemistry:12}
Differences of a few \unit{kJ/mol} involving a hydrogen bond need correlation and
dispersion: a hybrid functional with a dispersion correction, or better a
correlated wavefunction method for the final energies, with a polarised
triple-split basis including \omterm{def:b3:computational-chemistry:split-valence}{diffuse functions}. Optimise both conformers at
the same level, confirm minima by frequencies (all real), add zero-point
energies, test the basis by one larger calculation, and compare with a
related system of known conformational energy.
\end{solution}

\begin{solution}{pb:b3:computational-chemistry:1}
\textbf{1.} Each \omterm{def:b3:computational-chemistry:basis}{Slater-type orbital} of a minimal basis is replaced by a fixed
contraction of three Gaussians fitted to it.
\textbf{2.} The helium nucleus attracts its electrons more, so its $1s$ function
is more compact (larger exponents). $6.3624/3.4253 = 1.857 = (1.69/1.24)^2$.
\textbf{3.} Two basis functions, two electrons, one occupied orbital (and one
virtual).
\textbf{4.} $V_{\mathrm{nn}} = Z_{\mathrm{He}}Z_{\mathrm H}/R = 2/1.4632 =
\qty{1.3669}{\hartree}$.
\textbf{5.} The two functions overlap strongly (54\,\%): the atoms are close
enough to bond.
\textbf{6.} $h_{11}$ is the energy of an electron in the He function with both
nuclei but no other electron; the He nucleus (charge 2) holds it much more
tightly.
\textbf{7.} $(h_{11} - \varepsilon)(h_{22} - \varepsilon) - (h_{12} - \varepsilon S)^2 = 0$:
$0.71185\varepsilon^2 + 2.79292\varepsilon + 2.44969 = 0$.
\textbf{8.} $\varepsilon = \qty{-2.600}{\hartree}$ and \qty{-1.324}{\hartree}.
\textbf{9.} $\mathbf h$ ignores the repulsion between the two electrons; the core
orbital is too contracted and too low.
\textbf{10.} $F_{\mu\nu} = h_{\mu\nu} + \sum_{\lambda\sigma}P_{\lambda\sigma}[(\mu\nu|\sigma\lambda) -
\frac12(\mu\lambda|\sigma\nu)]$, with $P_{\lambda\sigma} = 2C_{\lambda1}C_{\sigma1}$.
\textbf{11.} The Coulomb repulsion of each electron by the charge cloud of the
other (and, in general, exchange).
\textbf{12.} $\mathbf F$ depends on $\mathbf P$, which depends on the orbitals obtained
from $\mathbf F$: the equations are non-linear and are solved by successive
approximations until self-consistency.
\textbf{13.} $c_1^2 + c_2^2 + 2c_1c_2S = 0.7684 + 0.0410 + 0.1906 = 1.0000$.
\textbf{14.} He: $1.5369 + 0.1906 = 1.727$; H: $0.0820 + 0.1906 = 0.273$
electron.
\textbf{15.} Charges: He $+0.27$, H $+0.73$. The positive charge sits mostly on
hydrogen: \ce{HeH+} is a proton bound to a helium atom, \ce{He + H+}, as expected
since the ionisation energy of He (24.6 eV) far exceeds that of H (13.6 eV).
\textbf{16.} $-\varepsilon_1 = \qty{1.633}{\hartree} = \qty{44.4}{eV}$.
\textbf{17.} $E_{\mathrm{el}} = -2.8418 - 1.3669 = \qty{-4.2087}{\hartree}$.
\textbf{18.} Three cycles reach \qty{-2.8418}{\hartree}. Each cycle gives a
determinant, whose energy is an upper bound to the converged Hartree--Fock
energy (\omterm{thm:b3:computational-chemistry:variational}{variational principle}), and the iterations improve it.
\textbf{19.} The basis with $\zeta = 2.0925$ gives the lower energy, so it is the
better one for this molecule (\omterm{thm:b3:computational-chemistry:variational}{variational principle}): the helium function
is more compact in \ce{HeH+} than in the free atom, whose $\zeta = 1.69$ the
standard basis copies.
\textbf{20.} The empty antibonding orbital $\sigma^*$; its energy would approximate
minus the electron affinity of \ce{HeH+} (poorly, in such a small basis).
\textbf{21.} Zero (no electron); $\qty{-2.8078}{\hartree}$.
\textbf{22.} $-2.8078 - (-2.8418) = \qty{0.0340}{\hartree} = \qty{89}{kJ/mol}$.
\textbf{23.} About half the measured value: a minimal basis cannot polarise the
helium density towards the proton (there are no $p$ functions), so the bond
is too weak. (Zero-point and \qty{298}{K} corrections are small beside this.)
\textbf{24.} $-(24.5874 + 54.4178)/27.211 = \qty{-2.9034}{\hartree}$; the STO-3G atom
is \qty{0.0956}{\hartree} (\qty{2.6}{eV}) too high.
\textbf{25.} Geometries depend on how the energy changes with $R$; most of the
error, concentrated in the core, is nearly the same at every geometry and
cancels.
\textbf{26.} \textbf{$E(\ce{HeH+}) = \qty{-2.8418}{\hartree}$ at RHF/STO-3G (standard
basis), \qty{1.4632}{\bohr}} (and \qty{-2.8607}{\hartree} with the classic
$\zeta_{\mathrm{He}} = 2.0925$).
\end{solution}
